% Options for packages loaded elsewhere
\PassOptionsToPackage{unicode}{hyperref}
\PassOptionsToPackage{hyphens}{url}
%
\documentclass[
  11pt,
  a4paper,
]{article}
\usepackage{amsmath,amssymb}
\usepackage{iftex}
\ifPDFTeX
  \usepackage[T1]{fontenc}
  \usepackage[utf8]{inputenc}
  \usepackage{textcomp} % provide euro and other symbols
\else % if luatex or xetex
  \usepackage{unicode-math} % this also loads fontspec
  \defaultfontfeatures{Scale=MatchLowercase}
  \defaultfontfeatures[\rmfamily]{Ligatures=TeX,Scale=1}
\fi
\usepackage{lmodern}
\ifPDFTeX\else
  % xetex/luatex font selection
\fi
% Use upquote if available, for straight quotes in verbatim environments
\IfFileExists{upquote.sty}{\usepackage{upquote}}{}
\IfFileExists{microtype.sty}{% use microtype if available
  \usepackage[]{microtype}
  \UseMicrotypeSet[protrusion]{basicmath} % disable protrusion for tt fonts
}{}
\makeatletter
\@ifundefined{KOMAClassName}{% if non-KOMA class
  \IfFileExists{parskip.sty}{%
    \usepackage{parskip}
  }{% else
    \setlength{\parindent}{0pt}
    \setlength{\parskip}{6pt plus 2pt minus 1pt}}
}{% if KOMA class
  \KOMAoptions{parskip=half}}
\makeatother
\usepackage{xcolor}
\usepackage[margin=26mm]{geometry}
\setlength{\emergencystretch}{3em} % prevent overfull lines
\providecommand{\tightlist}{%
  \setlength{\itemsep}{0pt}\setlength{\parskip}{0pt}}
\setcounter{secnumdepth}{5}
\usepackage{mathtools}
\usepackage{amssymb}
\usepackage{microtype}
\usepackage{fancyhdr}
\usepackage{longtable,booktabs,array}
\setlength{\LTpre}{8pt}
\setlength{\LTpost}{8pt}
\renewcommand{\arraystretch}{1.15}
\pagestyle{fancy}
\fancyhf{}
\fancyhead[L]{\footnotesize Paliy: Prime triangle-tiling counts}
\fancyhead[R]{\footnotesize v0.1.0 --- Draft for review}
\fancyfoot[C]{\thepage}
\renewcommand{\headrulewidth}{0.3pt}
\setlength{\headheight}{15pt}
\setlength{\emergencystretch}{2em}
\setcounter{tocdepth}{1}
\allowdisplaybreaks[1]
\AtBeginDocument{\hypersetup{urlcolor=blue,linkcolor=black,pdftitle={A candidate classification of prime triangle-tiling counts},pdfauthor={Denis Paliy}}}
\ifLuaTeX
  \usepackage{selnolig}  % disable illegal ligatures
\fi
\IfFileExists{bookmark.sty}{\usepackage{bookmark}}{\usepackage{hyperref}}
\IfFileExists{xurl.sty}{\usepackage{xurl}}{} % add URL line breaks if available
\urlstyle{same}
\hypersetup{
  pdftitle={A candidate classification of prime triangle-tiling counts},
  pdfauthor={Denis Paliy},
  hidelinks,
  pdfcreator={LaTeX via pandoc}}

\title{A candidate classification of prime triangle-tiling counts}
\author{Denis Paliy}
\date{29 September 2026}

\begin{document}
\maketitle

\begin{center}
\small\textsc{Version 0.1.0 --- Draft for review}
\end{center}

\begin{abstract}
We present a candidate proof that a prime $p$ is the number of congruent
triangles in a dissection of some triangle if and only if $p=2$, $p=3$,
or $p\equiv1\pmod4$. The two new proposed geometric inputs exclude the
primitive scale-one targets in the scalene $W$ and $\beta$-isosceles
branches for tiles satisfying $3\alpha+2\beta=\pi$. The isosceles
argument forces a standard corner whose complement is the forbidden
$W$ triangle; the $W$ argument uses a column induction. Both permit
arbitrary T-junctions. We give the complete passage from these inputs
and the published shape classification to the prime statement, including
explicit existence constructions and the arithmetic for the other
branches. The classification of all composite counts remains unresolved.
\end{abstract}

\noindent\textbf{Status and contribution.} This is a candidate proof,
released for mathematical scrutiny. Denis Paliy set the research
direction and guided the investigation; ChatGPT assisted with exploring
proof strategies, drafting arguments, and carrying out internal checks.
The draft has not received external peer review or formal verification.
No research-priority claim is made. The forced-patch and
column-completion arguments are the principal new steps requiring review.

\section{The prime case and its dependencies}\label{sec:prime}

A tiling is a finite dissection into triangles with pairwise disjoint
interiors, all congruent to one fixed tile. Reflections are allowed, and
an edge may meet vertices of several other tiles: no edge-to-edge
assumption is imposed.

\noindent\textbf{Proposed prime-classification theorem.}
For a prime $p$, there exists a triangle tiled by $p$ congruent triangles
if and only if
\[
\boxed{\quad p=2,\qquad p=3,\qquad\text{or}\qquad p\equiv1\pmod4.\quad}
\]

The proposed new ingredients are exactly the two scale-one exclusions
proved in Section~\ref{the-isosceles-obstruction} and
Appendix~\ref{complete-scale-one-w-exclusion}. This section explains all
outside dependencies and the full deduction from those ingredients. It
does not infer a classification of composite counts from the prime case.

\subsection{Published inputs and exhaustive branch list}

The shape classification is due to Laczkovich~\cite{laczkovich},
Theorems~4.1, 5.1 and 5.3. Its isosceles specialization is recorded in
\cite{isosceles}, Theorem~3.1 and Table~1; its non-isosceles
specialization is restated in \cite{prime}, Theorems~6 and~9, with an
independent comparison in \cite{rationality}, Table~1. The reptiling
classification is due to Snover, Waiveris and Williams~\cite{reptile};
see also \cite{isosceles}, Theorem~2.1.

We use the specified versions of the sources. In particular,
\cite{group1} means version~4: its corrected necessary equations are
used without assuming that a divisibility condition sufficient for a
particular construction is necessary for every tiling. Rationality in
the irrational $120^\circ$ branches is supplied by
Beeson--Zhang~\cite{rationality}, Theorem~1.1, rather than the flawed
older rationality argument.

Suppose $p>3$ is prime and a tiling exists. If the tile is similar to
the target, the reptiling classification gives a square, three times a
square, or a sum of two integer squares as the count. Only the last
option can equal $p>3$, and reduction modulo~4 then gives
$p\equiv1\pmod4$. For nonsimilar tilings the following list is
exhaustive. Here ``Group~1'' denotes $3\alpha+2\beta=\pi$.

\begin{longtable}{@{}>{\raggedright\arraybackslash}p{0.49\textwidth}>{\raggedright\arraybackslash}p{0.47\textwidth}@{}}
\toprule
\textbf{Target and tile branch} & \textbf{Exclusion for $p>3$}\\
\midrule
\endfirsthead
\toprule
\textbf{Target and tile branch (continued)} & \textbf{Exclusion for $p>3$}\\
\midrule
\endhead
\bottomrule
\endfoot
Equilateral; rational tile angles & Elementary area calculation below.\\[3pt]
Equilateral; irrational angles with a $60^\circ$ tile angle &
\cite{equilateral}, Theorem~3(i), and the factorization below.\\[3pt]
Equilateral; irrational angles with a $120^\circ$ tile angle &
\cite{equilateral}, Theorem~6, with rationality from \cite{rationality}.\\[3pt]
Non-equilateral isosceles; right tile &
\cite{isosceles}, Theorem~7.8: the count is square or even.\\[3pt]
Non-equilateral isosceles; tile $\gamma=2\alpha$ &
\cite{isosceles}, Theorem~11.7: the count is not squarefree.\\[3pt]
Group~1; target $(2\alpha,\alpha,2\beta)$ &
\cite{group1}, Theorem~16.\\[3pt]
Group~1; isosceles base angle $\alpha+\beta$ &
\cite{group1}, Theorem~22.\\[3pt]
Group~1; isosceles base angle $\alpha$ &
\cite{group1}, Theorem~26.\\[3pt]
Group~1; $W=(2\alpha,\beta,\alpha+\beta)$ &
Integer-scale reduction below, then Appendix~\ref{complete-scale-one-w-exclusion}.\\[3pt]
Group~1; isosceles base angle $\beta$ &
Integer-scale reduction below, then Section~\ref{the-isosceles-obstruction}.\\[3pt]
Irrational $120^\circ$ tile; non-equilateral isosceles target &
Signed-direction obstruction below.\\[3pt]
Irrational $120^\circ$ tile; non-isosceles target &
The four sine-law factorizations below.\\
\end{longtable}

To see exhaustiveness, first remove reptilings. The equilateral
classification gives precisely the five tile types covered by the first
three rows. For a non-equilateral isosceles target, the classification
gives right tiles, the $\gamma=2\alpha$ family, the three isosceles
Group~1 shapes, and $\gamma=120^\circ$. For a non-isosceles target with
rational multiples of $\pi$ as angles, the rational-angle theorem
permits only a reptiling. In the remaining case the tile also has an
irrational angle, because target angles are sums of tile angles; the
classification gives the two scalene Group~1 shapes and the four
scalene $120^\circ$ shapes. Labels may be permuted as allowed by these
classifications; the relation $3\alpha+2\beta=\pi$ is not silently
swapped.

\noindent\textbf{Source qualifications.} No broad prime-count conclusion
from \cite{prime} is used: its original Group~1 and isosceles
$120^\circ$ dependencies are insufficient. In particular, we do not use
its Theorem~3, Theorem~5(iii), Theorem~22 or Corollary~23 as an
established global result. The withdrawn paper arXiv:2607.19572 is not
used. Two smaller wording errors also require care: Corollary~7.9 of
\cite{isosceles} omits the valid $N=2$ exception, so we use
Theorem~7.8 for odd primes; and the rational-angle equilateral
discussion omits the $2k^2$ family of $30$--$60$--$90$ tiles, which is
included below.

\subsection{Direct integer-scale reduction in the two new branches}

Rationality and primitive normalization in Group~1 give
\[
(a,b,c)=(uv,v^2-u^2,v^2),\qquad
0<u<v,\quad\gcd(u,v)=1.
\]
Write $Q=2v^2-u^2$ and $P=3v^2-u^2$. The sine law gives the
side proportions
\[
W:\ (v^3,uQ,v(v^2-u^2)),\qquad
\beta\text{-isosceles}:\ (v^3,v^3,uP).
\]
Both triples are primitive: $\gcd(v^3,uQ)=\gcd(v^3,uP)=1$ follows
from $\gcd(u,v)=1$. Every actual target side is an integer sum of
integer tile-edge lengths. B\'ezout's identity therefore makes its
similarity scale a positive integer $j$. Comparison of areas yields
\[
N=j^2Q\quad\text{for }W,\qquad
N=j^2P\quad\text{for the }\beta\text{-isosceles shape}.
\]
This direct argument supplies the required integer scale independently
of coloring-divisibility claims. It agrees with the corrected formulas
in \cite{group1}, Corollary~1 and Theorem~19. Each of $P,Q$ exceeds
one, so a prime count forces $j=1$. The proposed geometric theorems
below exclude exactly these two targets at $j=1$, for every admissible
$u,v$. Neither uniqueness of representation by these quadratic forms
nor a construction-specific condition such as $v\mid u^2$ is needed.

\subsection{Classical equilateral cases and the irrational $60^\circ$ case}

For equilateral tiles the count is a square. For tile sides
$(1,1,\sqrt3)$, an equilateral target side has the form
$L=A+B\sqrt3$ with nonnegative integers $A,B$; its area ratio is
$N=L^2$. Rationality forces $AB=0$, giving a square or three times a
square. For tile sides $(1,\sqrt3,2)$, the ratio is $N=L^2/2$;
again $AB=0$, and integrality gives $N=2k^2$ or $6k^2$. None of
these counts is prime greater than three.

In the irrational $60^\circ$ branch, \cite{equilateral},
Theorem~3(i), gives integers $M,y$ satisfying
\[
y^2=(9p-M^2)(p-M^2),\qquad 0<M^2<p.
\]
Set $X=5p-M^2$. Then $(X-y)(X+y)=16p^2$, with both factors
positive. They cannot both be divisible by $p$, because their sum is
$10p-2M^2$ and $p\nmid M$. Thus one contains the full factor $p^2$,
so their sum is greater than $p^2$. But the sum is less than $10p$,
which is impossible for $p\ge11$. For $p=5$ the candidate
discriminants are $176,41$; for $p=7$ they are $372,177$. None is a
square. This checks every prime $p>3$ without relying on the printed
calculation in part~(ii) of that theorem.

\subsection{The irrational $120^\circ$ branches}

By rationality, normalize the tile to primitive positive integers
$a,b,c$ with
\[
c^2=a^2+ab+b^2.
\]
The three integers are pairwise coprime. Also $3\nmid c$: if
$3\mid c$, reduction modulo~3 gives $a\equiv b\pmod3$; writing
$a=b+3h$ yields $c^2=3b^2+9bh+9h^2$, which forces $3\mid b$
and contradicts primitiveness.

For the non-isosceles targets, the sine law gives the following
primitive integer side proportions and area ratios. Thus the actual
similarity scale $\lambda$ is a positive integer by the same boundary
length and B\'ezout argument as above.

\begin{center}
\small
\begin{tabular}{@{}lll@{}}
\toprule
\textbf{Target angles} & \textbf{Primitive side proportions} & $\boldsymbol{N}$\\
\midrule
$(\alpha,2\alpha,3\beta)$ & $(c^2,c(a+2b),3b(a+b))$ &
$3\lambda^2(a+2b)(a+b)$\\[3pt]
$(\alpha,2\beta,2\alpha+\beta)$ & $(ac,b(2a+b),(a+b)c)$ &
$\lambda^2(2a+b)(a+b)$\\[3pt]
$(\alpha,\alpha+\beta,\alpha+2\beta)$ & $(a,c,a+b)$ &
$\lambda^2(a+b)/b$\\[3pt]
$(2\alpha,2\beta,\alpha+\beta)$ & $(a(a+2b),b(2a+b),c^2)$ &
$\lambda^2(a+2b)(2a+b)$\\
\bottomrule
\end{tabular}
\end{center}

The first, second and fourth expressions are composite. In the third,
$b\mid\lambda^2$, and $a+b$ is composite: if $a+b=q$ were prime,
the cosine equation modulo $q$ would give $q\mid(c-a)(c+a)$.
Since $a<c<a+b=q$, the first factor cannot contain $q$; hence
$c+a=q$, giving $c=b$, a contradiction. These are the individual
computations of \cite{prime}, Theorems~18--21, without that paper's
global conclusion.

For a non-equilateral isosceles target, take the base angle to be
$\alpha$ after interchanging $a,b$ if necessary. Its primitive side
proportions are $(c,c,a+2b)$: any prime common to $c$ and $a+2b$
would divide $3b^2$, and $3\nmid c$. Thus its sides are
$k(c,c,a+2b)$ for an integer $k>0$, and
\[
N=\frac{k^2(a+2b)}{b}.
\]
Because $\gcd(b,a+2b)=1$, integrality and primality force
$k^2=b$ and $N=a+2b$.

We reproduce the signed-direction obstruction associated with
Bonfioli~\cite{bonfioli}. Fix a reference edge. Across a shared edge
segment, directions propagate through angle differences of the tile;
the adjacency graph of a finite dissection is connected. All directed
edge directions therefore lie, after an overall rotation, in
$G=\mathbb Z\pi/3+\mathbb Z\alpha$ modulo $2\pi$.
Irrationality of $\alpha/\pi$ makes the character
\[
f(j\pi/3+\ell\alpha)=(-1)^j
\]
well defined. Give an oriented segment of length $L$ the weight
$Lf(\text{direction})$. Reversal negates its weight. Subdivision at
every intersection and T-junction cancels all interior subsegments in
the sum of oriented tile boundaries. Each tile contributes
$\pm(c+a-b)$, while the target contributes $k(a+2b-2c)$.
Consequently
\[
\frac{k(a+2b-2c)}{c+a-b}
=\frac{c-a-b}{k}\in\mathbb Z,
\]
where the identity uses $k^2=b$ and the cosine equation. Hence
$k\mid(a+b-c)$. But $a<c<a+k^2$, so $c=a+kr$ for an integer
$1\le r\le k-1$. Substitution gives
\[
a(2r-k)=k(k-r)(k+r).
\]
Since $\gcd(a,k)=1$, we have $k\mid2r$. The only possible
multiple in $2\le2r\le2k-2$ is $2r=k$, making the left side
zero and the right side positive. For $k=1$ there is no possible $r$.
This excludes the isosceles branch, including T-junction tilings.

The equilateral irrational $120^\circ$ branch is excluded by
\cite{equilateral}, Theorem~6, with the rationality input stated
above. This completes all $120^\circ$ rows.

\subsection{Existence and completion of the conditional deduction}

For $p=2$, cut an isosceles triangle along its altitude. For $p=3$,
join the center of an equilateral triangle to its vertices; the pieces
are three congruent $120^\circ$ isosceles triangles.

For $p\equiv1\pmod4$, the sum-of-two-squares theorem gives
$p=m^2+n^2$ with positive integers $m,n$. Let the tile $T$ be the
right triangle with legs $m,n$. Take a larger right triangle $OAB$
with
\[
OA=m\sqrt p,\qquad OB=n\sqrt p,\qquad AB=p.
\]
If $H$ is the foot of the altitude from $O$ to $AB$, the projection
identities give
\[
AH=m^2,\qquad BH=n^2,\qquad OH=mn.
\]
Thus $OAH$ and $OBH$ are copies of $T$ enlarged by factors $m$
and $n$. Their standard parallel-line subdivisions contain $m^2$ and
$n^2$ triangles congruent to $T$, filling $OAB$ with exactly $p$
tiles.

Every nonsimilar branch has now either been excluded by a stated
published input and the calculations above, or reduced to one of the
two proposed geometric exclusions below. With those exclusions, a
prime $p>3$ allowing a tiling must be a reptiling count and satisfy
$p\equiv1\pmod4$. The constructions prove the converse. The remaining
mathematical burden of this draft is therefore concentrated in the two
scale-one proofs that follow.

\hypertarget{the-isosceles-obstruction}{%
\section{The isosceles obstruction}\label{the-isosceles-obstruction}}

\hypertarget{main-theorem-and-its-geometric-idea}{%
\subsection{Main theorem and its geometric
idea}\label{main-theorem-and-its-geometric-idea}}

Let \(0<u<v\), \(\gcd(u,v)=1\), and put

\[
a=uv,\qquad b=v^2-u^2,\qquad c=v^2,\qquad
P=3v^2-u^2,\qquad Q=2v^2-u^2.
\]

Let the congruent tile have sides \((a,b,c)\) opposite angles
\((\alpha,\beta,\gamma)\). Then \(3\alpha+2\beta=\pi\) and
\(\gamma=2\alpha+\beta>\pi/2\).

\textbf{Theorem.} The isosceles triangle with sides \((v^3,v^3,uP)\)
cannot be tiled by these tiles. Its area would require exactly \(P\)
tiles. The triangle with sides \((v^3,uQ,vb)\) cannot be tiled by these
tiles either; its area would require \(Q\) tiles.

The second assertion is the scale-one W theorem, proved in the appendix.
The argument below reduces the first assertion to the second by an
actual dissection: starting at the three-\(\alpha\) apex, it forces a
standard \(v^2\)-tile corner. Its complement is exactly the forbidden W
triangle.

These are statements about the full scale-one families. They do not
classify larger-scale composite tilings or by themselves solve every
part of Erdős problem 634.

\hypertarget{elementary-angular-and-arithmetic-facts}{%
\subsection{Elementary angular and arithmetic
facts}\label{elementary-angular-and-arithmetic-facts}}

The number \(\alpha/\pi\) is irrational: \(2\cos\alpha=2-u^2/v^2\) is
rational strictly between 1 and 2, whereas for a rational multiple of
\(\pi\) it would be a rational algebraic integer, hence an integer.
Consequently the tile-angle inventories at a \(\beta\) corner, an
\(\alpha\) sector, a \(\theta=\alpha+\beta\) sector, and a \(3\alpha\)
apex are respectively one \(\beta\), one \(\alpha\), one \(\alpha\) plus
one \(\beta\), and three \(\alpha\) sectors.

We need three exact positive-chain facts. All coefficients count whole
edges and are nonnegative integers.

\[
ib=Aa+Bb+Cc,\quad0<i<v
\ \Longrightarrow\ A=C=0,\ B=i. \tag{1}
\]

Indeed \(B\equiv i\pmod v\) and \(0\le B\le i\).

\[
ia=Aa+Bb+Cc,\quad0<i<v
\ \Longrightarrow\ B=C=0,\ A=i. \tag{2}
\]

Write \(B=v\ell\) and \(\ell v+C=T u\); division of the length equation
gives \(i=A+(T-\ell)v+\ell(v-u)\). If \(\ell+C>0\), then
\(Tu=\ell v+C>\ell u\), so \(T>\ell\) and \(i\ge v\), impossible.

\[
\begin{gathered}nc=Aa+Bb+Cc,\quad0<n<v\\
\Longrightarrow\quad B=0,\quad A=vh,\quad C=n-uh\quad\text{for an integer }h\ge0.\end{gathered}\tag{3}
\]

To prove (3), write \(B=v\ell\) and \(d=v-u\). The equation gives

\[
A=vh-d\ell,\qquad n=C+uh+d\ell.
\]

If \(\ell>0\), nonnegativity of \(A\) forces \(h\ge1\), and then
\(n\ge u+d=v\). Thus \(\ell=0\). In particular, when \(n<u\), (3) also
forces \(A=0\) and \(C=n\).

The distinction is crucial: below height \(u\) a pure-\(c\) seam has
only \(c\) opposite it; below height \(v\) it may have \(a\) opposite
it, but it has \textbf{no \(b\) opposite it}. The downward apex
induction needs the latter, stronger range.

Throughout the proof the tiling is fixed. A collinear component means a
connected union of actual whole tile edges on one supporting line. An
internal maximal component has a complete chain on each side: an edge
cannot end inside an edge on the other side without another collinear
edge continuing the remaining straight sector. Nor can a supplied next
edge be swallowed by a tile at a junction. Maximality prevents an
endpoint from lying inside any constituent edge. When the proof instead
uses a subsegment, it separately establishes endpoints for both chains,
using external supporting lines or the starting edge of an actual tile.
Thus none of the length lemmas is applied to a chain of merely partial
edges.

\hypertarget{equal-side-boundary-fact-and-coordinates}{%
\subsection{Equal-side boundary fact and
coordinates}\label{equal-side-boundary-fact-and-coordinates}}

No corner of the isosceles target contains a \(\gamma\) sector: its two
base corners contain one \(\beta\) each, and its apex contains exactly
three \(\alpha\) sectors. This remains true when the total apex angle
\(3\alpha\) exceeds \(\gamma\).

Every \(a\)- or \(b\)-edge on a target side contributes a \(\gamma\)
endpoint, and two such endpoints cannot coincide at a straight boundary
junction because \(2\gamma>\pi\). A side consisting of \(n\) edges has
only \(n-1\) internal junctions. Hence every side contains at least one
\(c\)-edge.

An equal side has length \(v^3=vc\). Reduction modulo \(v\) makes its
\(b\)-count a multiple of \(v\). If it had at least \(v\) \(b\)-edges,
removing \(v\) of them would leave a positive whole-edge decomposition
of

\[
vc-vb=u a.
\]

Since \(u<v\), (2) makes every remaining edge an \(a\)-edge,
contradicting the presence of a \(c\)-edge. Therefore \textbf{both equal
sides have no \(b\)-edge}.

Place the target at \(A=0\), \(B=vq\), \(C=uP\), where \(p=a\) is
horizontal, \(q=ce^{i\beta}\), and \(q-p=be^{i\theta}\). Put
\(Z(i,t)=ip+tq\). Use

\[
D(i,t)=\operatorname{conv}(Z(i,t),Z(i+1,t),Z(i,t+1)),
\] \[
U(i,t)=\operatorname{conv}(Z(i+1,t),Z(i,t+1),Z(i+1,t+1)).
\]

Their angles at the listed vertices are respectively
\((\beta,\gamma,\alpha)\) and \((\alpha,\gamma,\beta)\).

The apex has exactly three \(\alpha\) tiles. Its two outer edges are
\(c\), because the equal sides have no \(b\). Each \(\alpha\) tile has
incident sides \(b,c\), so its two inner radial seams comprise one
\(b/b\) seam and one \(b/c\) seam. Choose the reflected labeling so the
latter is next to the left outer tile \(D(0,v-1)\).

\hypertarget{reverse-apex-state}{%
\subsection{Reverse-apex state}\label{reverse-apex-state}}

Write \(V_m=Z(m,v-m)\) for \(0\le m\le v\). Let \(T(m,t)\) be the actual
standard truncated corner region consisting of:

\begin{itemize}
\tightlist
\item
  all \(D(i,s)\) with \(0\le i<m\), \(s\ge t\), and \(i+s\le v-1\);
\item
  all \(U(i,s)\) with \(0\le i<m\), \(s\ge t\), and \(i+s\le v-2\).
\end{itemize}

Geometrically this is the standard subdivision of the polygon with
oblique coordinates

\[
(0,t),\quad(m,t),\quad(m,v-m),\quad(0,v),
\]

when \(t\le v-m\), with zero-height edges allowed.

For each \(m=1,\ldots,v-1\), the apex diagonal \(b\)-chain
\(B=V_0,V_1,\ldots,V_m\) will be actual. Its opposite chain starts at
\(B\) with the central apex \(c\)-edge. Since \(m<v\), short \(b\)
purity implies that \(V_m\) is in the relative interior of an actual
opposite edge. Thus the positive \(q\)-ray from \(V_m\) enters that
opposite tile's interior.

Let \(M_m\) be the actual maximal collinear component from \(V_m\) in
direction \(-q\). Its upper end is \(V_m\), by the preceding crossing.
We prove:

\begin{enumerate}
\def\labelenumi{\arabic{enumi}.}
\tightlist
\item
  Its whole left side is a chain of normal \(c\)-edges of \(U\) tiles.
\item
  If its lower endpoint is \(Z(m,r)\), then \(T(m,r)\) is actual; more
  generally every reached integer height \(t\) supplies \(T(m,t)\).
\item
  Its opposite side has no \(b\)-edge.
\end{enumerate}

These statements are proved in increasing \(m\); within each column, the
continuation goes downward. Earlier columns may terminate earlier, but
whenever a newly supplied \(U\) extends one of them downward, its
already proved conditional completion statement applies. No assumption
is made that any column reaches the base.

\hypertarget{reverse-horizontal-cap}{%
\subsection{Reverse horizontal cap}\label{reverse-horizontal-cap}}

The first outer apex tile is \(D(0,v-1)\). Its \(b\)-edge is \(BV_1\),
and the opposite central \(c\)-edge starts at \(B\) and overruns \(V_1\)
since \(c>b\).

More generally, suppose the diagonal has just reached \(V_m\), where
\(m<v\), and \(T(m-1,v-m)\) is known. For \(m=1\), the old region
\(T(0,v-1)\) is empty and the new tile \(D(0,v-1)\) is the outer apex
tile. The new \(D(m-1,v-m)\) gives, together with the old region, a
horizontal chain of \(m\) \(a\)-edges from \((v-m)q\) to \(V_m\).

Its negative extension at the outer side leaves the convex target. Its
positive extension at \(V_m\) enters the crossing opposite tile. Thus
the opposite chain is a complete chain of total length \(ma\), hence
exactly \(m\) \(a\)-edges.

At \(V_m\) the already placed \(D\) contributes \(\gamma\) and the
crossing tile a straight sector \(\pi\), so the remaining sector is
\(\theta=\alpha+\beta<\gamma\). Consequently the opposite \(a\)-tile has
\(\beta\), not \(\gamma\), at \(V_m\); it is \(U(m-1,v-m-1)\).
Propagating left, every successive opposite \(a\)-tile must also be a
normal \(U\): a reversed one would put two \(\gamma\) sectors in the
same lower half-plane at a shared junction.

At the outer boundary point \((v-m)q\), the old \(D\) contributes
\(\beta\) and the new \(U\) contributes \(\gamma\), leaving \(\alpha\).
The next boundary edge is \(c\), since that equal side has no \(b\);
this gives \(D(0,v-m-1)\).

At each interior grid point \(Z(i,v-m)\), where \(1\le i<m\), the
existing upper row and the two new \(U\) tiles leave exactly one
\(\alpha\) sector. The new left \(U\) supplies an actual downward
\(c\)-edge in the earlier column \(M_i\), connected to its upper
component through the old actual region. Its already proved
no-opposite-\(b\) property excludes \(b\) on the right of that
\(q\)-edge. Hence its \(\alpha\) tile uses \(c\) and is \(D(i,v-m-1)\).

Explicitly, at a reached height \(t\), the old region contains the
collinear \(c\)-edges of \(U(i-1,s)\) for \(t\le s\le v-i-1\), joining
\(Z(i,t)\) to \(V_i\). The new edge of \(U(i-1,t-1)\) attaches to this
very component. If it extended below an earlier claimed endpoint, that
would contradict the already established maximality; it would not create
a different column. The same connectivity applies in the downward
continuation below.

This supplies \(T(m,v-m-1)\) and the first complete \(c\)-edge of
\(M_m\).

\hypertarget{downward-column-completion}{%
\subsection{Downward column
completion}\label{downward-column-completion}}

Assume \(T(m,t)\) is actual, with \(t\le v-m-1\), and the real component
\(M_m\) continues downward from \(X=Z(m,t)\). Its old left-side normal
tiles \(D(m-1,t)\) and \(U(m-1,t)\) contribute
\(\gamma+\alpha=\pi-\beta\). The actual downward \(q\)-ray therefore
leaves one \(\beta\) tile on its left.

That \(\beta\) tile has either \(a\) or \(c\) on the downward \(q\)-ray.
If it had \(a\), its other incident side would be a \(c\)-edge from
\(X\) horizontally to the left. But the actual upper horizontal chain
from \(tq\) to \(X\) consists of \(m\) \(a\)-edges. Starting at \(X\)
with this \(c\)-edge, the opposite chain must cover the whole upper
chain: at a T-junction inside an upper edge a straight sector forces
continuation, and an upper-chain grid endpoint cannot let the actual
next upper edge be swallowed by a tile. At \(tq\) the external side
prevents overrun. Hence the lower chain would be a complete positive
chain of length \(ma\) containing a \(c\)-edge, contrary to short \(a\)
purity, since \(m<v\).

Thus the \(\beta\) tile has \(c\) on the \(q\)-ray and is
\(U(m-1,t-1)\). Its top \(a\)-edge starts a whole lower chain along the
existing upper chain of length \(ma\). Short \(a\) purity makes the
lower chain exactly \(m\) \(a\)-edges. Starting at \(X\), the same
two-\(\gamma\) argument forces all these lower tiles to be normal
\(U(i,t-1)\).

At \(tq\), the old \(D(0,t)\) contributes \(\beta\) and \(U(0,t-1)\)
contributes \(\gamma\), leaving \(\alpha\), so absence of boundary \(b\)
gives \(D(0,t-1)\). At each inner point \(Z(i,t)\), where \(1\le i<m\),
exactly one \(\alpha\) remains. The new actual downward \(c\)-edge is on
the earlier column \(M_i\); its no-opposite-\(b\) property makes that
\(\alpha\) tile \(D(i,t-1)\). Therefore \(T(m,t-1)\) is actual.

Every actual continuation adds one complete normal \(c\)-edge. Since the
target is bounded, the component ends at some grid height \(r\ge0\). Its
length is \(nc\), where

\[
n=(v-m)-r,\qquad 0<n\le v-m<v.
\]

Arithmetic fact (3) gives no \(b\)-edge on the opposite side. This
closes the induction for \(M_m\).

\hypertarget{advancing-the-apex-diagonal}{%
\subsection{Advancing the apex
diagonal}\label{advancing-the-apex-diagonal}}

At \(V_m\) the crossing opposite tile, the old \(D(m-1,v-m)\), and the
first \(U(m-1,v-m-1)\) leave one \(\alpha\) sector immediately on the
right of the downward \(q\)-ray. Its \(q\)-side cannot be \(b\) by
\(M_m\)'s just established no-opposite-\(b\) property. It is therefore
\(c\), and the tile is \(D(m,v-m-1)\). Its \(b\)-edge extends the actual
diagonal to \(V_{m+1}\).

For \(m+1<v\), short \(b\) purity again says that this endpoint lies
inside an actual opposite edge. The reverse horizontal cap then
initializes the next column. For \(m=v-1\) the same forcing gives
\(D(v-1,0)\) and reaches \(V_v=vp\) on the base; no short \(b\)
statement with index \(v\) is used.

The column \(M_{v-1}\) starts at height \(1\) and its initialization
reaches height \(0\). Its completion statement gives \(T(v-1,0)\).
Adding \(D(v-1,0)\) gives the entire standard

\[
K_v=\operatorname{conv}(0,vp,vq),
\]

with its \(v^2\) tiles.

In particular the base begins at \(A\) with \(v\) consecutive
\(a\)-edges, of total length \(L=va=uc\).

\hypertarget{the-complement-is-a-forbidden-w-triangle}{%
\subsection{The complement is a forbidden W
triangle}\label{the-complement-is-a-forbidden-w-triangle}}

Let \(L=vp\). The forced \(K_v\) has vertices \(A,B,L\), consists of
\(v^2\) whole tiles, and is a genuine corner of the target. Its diagonal
\(BL\) is the actual chain of \(v\) \(b\)-edges. Thus removing its tiles
leaves a tiling of the triangle \(BLC\).

But

\[
BC=v^3,\qquad BL=vb,\qquad
LC=uP-va=u(2v^2-u^2)=uQ.
\]

The remaining tile count is \(P-v^2=Q\). Hence the complement is exactly
the scale-one W triangle excluded below. This proves the isosceles
assertion. No base-edge count, two-\(c\) boundary lemma or assumption
about an ordering of the base edges is needed. \(\square\)

\clearpage
\appendix

\hypertarget{complete-scale-one-w-exclusion}{%
\section{Complete scale-one W
exclusion}\label{complete-scale-one-w-exclusion}}

The following proof is reproduced to make the result self-contained. It
uses a different orientation of coordinates, reset in its statement. The
arithmetic facts called \(S_b,S_a,S_c\) below are respectively (1), (2),
and the \(n<u\) consequence of (3).

\hypertarget{statement-and-coordinate-convention}{%
\subsection{Statement and coordinate
convention}\label{statement-and-coordinate-convention}}

Let \(0<u<v\), \(\gcd(u,v)=1\), and put

\[
a=uv,\quad b=v^2-u^2,\quad c=v^2,\quad Q=2v^2-u^2.
\]

The theorem is that the triangle with side lengths

\[
(v^3,\;uQ,\;vb)
\]

cannot be tiled by \(Q\) congruent copies of \((a,b,c)\). Its angles are
\((\beta,2\alpha,\alpha+\beta)\), where \(3\alpha+2\beta=\pi\).

Name the \(\beta\) corner \(O\). Direct \(OA\) horizontally to the
right, with \(|OA|=v^3\). The other side \(OC\) has length \(u(b+c)=uQ\)
and makes angle \(\beta\) with \(OA\). Put

\[
p=a,\quad q=ce^{i\beta},\quad O=0,\quad A=v^3,\quad
C=\frac{u(b+c)}{c}q,\quad Z_{i,t}=ip+tq.
\]

Thus \(q-p=be^{i\theta}\), where \(\theta=\alpha+\beta\). The \(q\)
direction is along the side \(OC\), whose edge counts will be
\((0,u,u)\). It is \textbf{not} assumed to be a pure \(c\)-grid
initially.

The two adjacent side lengths satisfy the strict inequalities

\[
\begin{aligned}|OA|&=v^3>u a=u^2v,\\
|OC|&=u(b+c)>u c.\end{aligned}\tag{L}
\]

Consequently every corner front of index \(j\le u\) has both endpoints
strictly inside the two external sides, so its opposite side is an
internal chord. This excludes the degenerate situation in which a
purported front is already the outer third side of the target.

Use the standard actual triangles

\[
D_{i,t}=\operatorname{conv}(Z_{i,t},Z_{i+1,t},Z_{i,t+1}),\quad
U_{i,t}=\operatorname{conv}(Z_{i+1,t},Z_{i,t+1},Z_{i+1,t+1}).
\]

The angles of \(D\) at its listed vertices are
\((\beta,\gamma,\alpha)\); those of \(U\) are \((\alpha,\gamma,\beta)\).
Let \(R(m,n)\) denote the actual rectangle of all such tiles for
\(0\le i<m\), \(0\le t<n\).

\hypertarget{boundary-counts-corner-and-short-chain-arithmetic}{%
\subsection{Boundary counts, corner, and short-chain
arithmetic}\label{boundary-counts-corner-and-short-chain-arithmetic}}

Every target angle is less than \(\gamma\). Thus no boundary corner can
contain a \(\gamma\) sector. Each boundary \(a\)- or \(b\)-edge
contributes a \(\gamma\) endpoint. Two \(\gamma\) sectors cannot share a
straight boundary junction. Consequently every target side has at least
one \(c\)-edge.

On \(OC\), write its counts as \((A_0,B_0,C_0)\), with \(C_0\ge1\). The
length equation is

\[
A_0uv+B_0(v^2-u^2)+C_0v^2=u(2v^2-u^2).
\]

Reduction modulo \(v\) gives \(B_0=u+vk\), \(k\ge0\). Write
\(C_0=u\ell-vk\); then \(A_0=uk+v(1-\ell)\). For \(k\ge1\), \(C_0\ge1\)
requires \(\ell\ge k+1\), but \(A_0\ge0\) requires \(\ell\le k\). Hence
\(k=0\), and positivity gives \(\ell=1\). Therefore

\[
\boxed{(A_0,B_0,C_0)=(0,u,u).}
\tag{B}
\]

At \(O\), the unique \(\beta\) tile has incident sides \(a,c\). Since
\(OC\) has no \(a\)-edge, its first edge is \(c\) and \(OA\) starts with
\(a\). Thus \(D_{0,0}\) is present.

We use \(S_b\), \(S_a\) and \(S_c\) for the positive-chain facts (1),
(2), and the \(n<u\) consequence of (3), proved above. The exact angle
inventories and irrationality of \(\alpha/\pi\) were also established
above.

\hypertarget{first-column-actual-c-only-strip-then-no-right-a-edge}{%
\subsection{First column: actual c-only strip, then no right
a-edge}\label{first-column-actual-c-only-strip-then-no-right-a-edge}}

The internal \(b\)-edge of \(D_{0,0}\) is the chord joining \(p\) and
\(q\). Both endpoints are on the outer boundary. Its opposite chain
consists of a single \(b\)-edge by \(S_b\). The opposite tile is
\(U_{0,0}\): the other orientation would put two \(\gamma\) sectors at
the boundary point \(p\).

Let \(L_1\) be the actual maximal connected collinear component of tile
boundaries from \(p\) in direction \(q\). Its negative extension leaves
the target across \(OA\). Its left side starts with the \(c\)-edge of
\(U_{0,0}\).

The distance from this line to the supporting line \(OC\) is
\(a\sin\beta=h_c\), the tile altitude to \(c\). A tile attached on the
left with an \(a\)- or \(b\)-edge would have larger altitude and cross
the external supporting line. Hence \textbf{every complete edge on the
entire left of \(L_1\) has length \(c\)}, even at arbitrary T-junctions.

The component is finite. Maximality means its left edges form a complete
chain starting at \(p\) and ending at \(p+nq\), for some integer
\(n\ge1\). There is no possibility of ending inside one of these edges,
because that edge would extend the component.

For each left \(c\)-edge from \(p+iq\) to \(p+(i+1)q\), its opposite
vertex lies on the supporting line \(OC\), since its altitude equals the
strip width. The two possible vertex positions are

\[
(i+1)q,\qquad(i+\lambda)q,
\qquad \lambda=2\frac ac\cos\beta
=\frac{u^2(3v^2-u^2)}{v^4}.
\]

The second point has signed distance

\[
(i+\lambda)c=iv^2+\frac{u^2(3v^2-u^2)}{v^2}
\]

from \(O\), which is nonintegral because \(\gcd(u,v)=1\). An actual tile
vertex on \(OC\) must be a boundary junction, and every
boundary-junction distance is an integer sum of tile-edge lengths. Thus
the second point is impossible. All left tiles are the normal
\(U_{0,i}\).

The spaces between consecutive \(U\) tiles and \(OC\) are the triangles
\(D_{0,i}\). They are bounded by actual tile edges and the external
boundary and each has exactly one tile's area. They must therefore be
single tiles, so \(R(1,n)\) is present. Equivalently, this can be
checked by successive boundary angle fans. The original \(D_{0,0}\) is
already present.

At \(nq\) the tiles \(D_{0,n-1}\) and \(U_{0,n-1}\) contribute
\(\alpha\) and \(\gamma\), leaving \(\beta\) in the target boundary
half-plane. This point cannot be \(C\): its two existing angles already
total \(\alpha+\gamma>\alpha+\beta\), the target angle at \(C\). It
cannot lie beyond \(C\) either. Hence it is an internal boundary
junction, and its next boundary edge belongs to a \(\beta\) tile. That
edge is \(a\) or \(c\); by (B), it is \(c\). Therefore the existing
first \(n\) boundary \(c\)-edges require one additional \(c\)-edge,
giving

\[
\boxed{n<u.}
\tag{H}
\]

The opposite side of \(L_1\) is a complete edge chain of length \(nc\).
By \(S_c\), every edge on it is also \(c\). In particular, there is
\textbf{no right \(a\)-edge on \(L_1\)}.

The same argument gives the first-column rectangular completion
property: whenever this actual component reaches height \(t\),
\(R(1,t)\) is present. If \(u=1\), (H) contradicts \(n\ge1\) already;
hence that case is excluded at this point.

\hypertarget{horizontal-cap-without-a-pre-assumed-c-grid}{%
\subsection{Horizontal cap without a pre-assumed
c-grid}\label{horizontal-cap-without-a-pre-assumed-c-grid}}

Suppose \(R(i,t)\) and \(D_{i,t-1}\) are present, \(0<i<v\), and an
actual \(\theta\)-directed edge contains \(X=Z_{i,t}\) in its relative
interior, with its tile on the lower/right side.

The rectangle supplies a horizontal chain of \(i\) lower \(a\)-edges
from \(tq\) to \(X\). The negative extension at \(tq\) leaves the target
across \(OC\), and the positive extension at \(X\) enters the crossing
tile's interior. Hence this is the entire actual horizontal component.
Its complete upper chain has length \(ia\), so \(S_a\) makes it exactly
\(i\) \(a\)-edges.

At \(tq\) the existing \(\alpha\) and \(\gamma\) sectors leave
\(\beta\). The point is an internal point of \(OC\) for the same reason
as in (H). Since \(OC\) has no \(a\)-edges, its next actual edge must be
\(c\), giving \(D_{0,t}\). At every later upper-chain junction, the
preceding normal \(D\) contributes \(\gamma\). The remaining upper angle
is \(\theta\), and \(\gamma>\theta\). The next upper \(a\)-tile must
therefore have \(\beta\) at that junction, so it is normal. This forces
\(D_{i-1,t}\).

Thus every needed local \(c\)-boundary continuation is forced by the
actual rectangle and (B), rather than assumed globally.

\hypertarget{diagonal-chain-obstruction}{%
\subsection{Diagonal-chain
obstruction}\label{diagonal-chain-obstruction}}

\textbf{Diagonal-chain lemma.} Assume \(1\le j<v\), the actual
\(R(j-1,n+1)\) and \(D_{j-1,n}\) are present, and every column with
index \(<j\) already has the following conditional rectangular
completion property: whenever its actual component connected to the base
reaches \(Z_{i,t}\), the whole \(R(i,t)\) is present. There cannot be an
actual \(a\)- or \(c\)-edge from \(V=Z_{j,n}\) in direction
\(e^{i\theta}\) toward \(OC\), with its tile on the lower/right side.

\textbf{Proof.} Set

\[
X_k=Z_{j-k,n+k},\qquad0\le k\le j.
\]

The supplied \(D_{j-1,n}\) gives the first actual \(b\)-edge \(X_0X_1\)
on the upper/left side. The hypothesized \(a\)- or \(c\)-edge starts at
\(X_0\) on its other side.

Suppose the actual upper/left \(b\)-chain \(X_0\cdots X_k\) has been
supplied. Along each of its edges, the tiling on the opposite side
supplies a collinear chain. If an opposite edge ends inside one of the
existing \(b\)-edges, that \(b\)-tile has a straight sector there and
forces the next actual collinear edge. Hence the opposite chain starts
at \(X_0\) and covers the entire segment \(X_0X_k\).

If \(X_k\) were the endpoint of an opposite edge too, both sides from
\(X_0\) to \(X_k\) would be complete whole-edge chains of length \(kb\).
The opposite chain contains its initial \(a\)- or \(c\)-edge; because
\(k\le j<v\), this contradicts \(S_b\). Therefore every intermediate
\(X_k\) lies strictly inside one actual opposite edge.

At \(k=1\), put \(i=j-1\) and \(t=n+1\). The rectangle and additional
tile required by the horizontal-cap lemma are exactly the supplied
\(R(j-1,n+1)\) and \(D_{j-1,n}\). The lemma forces \(D_{i-1,t}\), whose
\(b\)-edge is \(X_kX_{k+1}\).

For \(i-1>0\), this newly forced tile also supplies a real \(c\)-edge in
column \(i-1\) from height \(t\) to \(t+1\). Its lower endpoint is
already connected to the base through \(R(i,t)\). The completion
property of this strictly smaller column therefore supplies
\(R(i-1,t+1)\). The new \(D_{i-1,t}\) is precisely the additional tile
needed at the next crossing. Thus the cap can be applied successively.
If any required tile leaves the target, this already contradicts the
assumed tiling.

When \(k=j\), the endpoint \(X_j=Z_{0,n+j}\) lies on the supporting line
\(OC\). If it were beyond the side segment, an already supplied tile
would have left the target. Otherwise the actual \(b\)-chain reaches the
external side there. Its opposite chain cannot overrun outside the
convex target, so both chains are complete from \(X_0\) to \(X_j\).
Their length \(jb\) contradicts \(S_b\) once more. The case \(j=1\) uses
this final step immediately, without an intermediate cap. \(\square\)

This proof permits intermediate endpoints and changes of the actual
opposite tile. It requires neither \(a>b\), nor a short-distance
estimate, nor artificial continuation of the original \(a\)- or
\(c\)-edge.

\hypertarget{strong-column-induction}{%
\subsection{Strong column induction}\label{strong-column-induction}}

Suppose \(K_m=\operatorname{conv}(0,mp,mq)\), with its standard
subdivision, is present and \(m\le u<v\). By (L), its front from \(mp\)
to \(mq\) is an internal chord. Its complete opposite edge chain is
exactly \(m\) \(b\)-edges by \(S_b\). At \(mp\) the old patch
contributes \(\gamma\), so the first opposite tile must have \(\alpha\)
there; another \(\gamma\) would exceed the boundary angle \(\pi\). Reversing
the orientation of a later opposite tile would put two \(\gamma\)
sectors in the same half-plane at a shared endpoint. Hence every
opposite tile is the normal \(U\) tile. This complete opposite row,
together with \(K_m\), supplies \(R(m,1)\).

For every actual column \(L_m\) from \(mp\) in direction \(q\),
establish:

\begin{enumerate}
\def\labelenumi{\arabic{enumi}.}
\tightlist
\item
  its whole left side consists of normal \(U\) tiles with complete
  \(c\)-edges;
\item
  if its connected prefix reaches height \(t\), \(R(m,t)\) is present;
\item
  it has no right \(a\)-edge.
\end{enumerate}

Column one has already been proved. Assume every smaller column is
established, and maintain an actual \(R(m,n)\). If \(L_m\) ends at
\(Z_{m,n}\), its completed normal left chain stops. Otherwise its actual
upward \(q\)-ray exists.

At \(nq\), the boundary \(\alpha\) and \(\gamma\) sectors force a normal
\(\beta\) tile \(D_{0,n}\), using (B). Sweep across the top row. At each
\(Z_{i,n}\), \(1\le i<m\), the old rectangle and preceding new \(D\)
leave the fan \(\theta\). In the following two cases the fan is read
counterclockwise from the positive horizontal ray to the \(\theta\) ray.

\begin{itemize}
\tightlist
\item
  The order \(\alpha\),\(\beta\) gives a real long \(a\)- or \(c\)-edge
  in the \(\theta\) direction, on its lower/right side. The
  diagonal-chain obstruction applies with \(j=i\); the preceding sweep
  has supplied its required rectangle and \(D\) tile. Thus this order is
  impossible.
\item
  In the order \(\beta\),\(\alpha\), an actual upward \(q\)-ray is
  present and connected to column \(i\) through the old rectangle. The
  \(\beta\) tile cannot use \(a\) on that ray, by the no-right-\(a\)
  property of this earlier column. Hence it is normal \(D_{i,n}\). The
  \(\alpha\) tile cannot use \(b\) on the ray, by that earlier column's
  complete c-only left side. Hence it is normal \(U_{i-1,n}\).
\end{itemize}

At the outer point \(Z_{m,n}\), its actual continuation leaves an
\(\alpha\) fan on the left. If the \(\alpha\) tile uses \(b\) on the
q-ray, its \(c\)-edge is a real forbidden diagonal; the just completed
sweep supplies the hypotheses of the diagonal-chain obstruction with
\(j=m\). Thus it uses \(c\), and is \(U_{m-1,n}\). This gives
\(R(m,n+1)\).

Every real continuation therefore adds a whole normal \(c\)-edge. Since
the target is bounded, continuation is finite and the full component has
some integer height \(n\) and the complete actual rectangle \(R(m,n)\).
Its boundary endpoint \(nq\) again forces one additional boundary
\(c\)-edge by (B), so \(n<u\). The right chain of the component has
length \(nc\); \(S_c\) makes it c-only. In particular it contains no
right \(a\)-edge. This completes the induction using only earlier
columns.

\hypertarget{forced-fronts-and-contradiction}{%
\subsection{Forced fronts and
contradiction}\label{forced-fronts-and-contradiction}}

Start with \(K_1=D_{0,0}\). Across its b-chord, the opposite tile is
forced. The \(\beta\) gap at \(q\) has its boundary edge on \(OC\) and
must use \(c\) there by (B). The \(\beta\) gap at \(p\) cannot switch to
a right a-edge on \(L_1\). Hence \(K_2\) is forced whenever \(u\ge2\).

Suppose \(K_j\) is present for \(j\le u<v\). By (L), its two front
endpoints are strictly inside the external sides and its inner side is
an actual internal chord. The complete opposite b-row is forced by
\(S_b\). Each remaining \(\beta\) notch admits the standard order of
sides or an interchange of \(a,c\). The notch on \(OC\) must be standard
because that side has no \(a\)-edge. Any other switched notch gives a
right a-edge on the actual column connected to the base through \(K_j\)
and its real opposite b-row. Its column index is at most \(j\), and all
such columns have just been established from their contained \(K_m\).
Therefore no switch is possible and \(K_{j+1}\) is forced.

Thus \(K_{u+1}\) is forced. It supplies \(u+1\) complete c-edges on
\(OC\), contradicting its exact count \(u\) from (B). This concludes the
proof.

\hypertarget{scope-and-relation-to-parametrization}{%
\section*{Scope and relation to parametrization}\label{scope-and-relation-to-parametrization}}

The two obstructions apply to the full primitive scale-one families
stated above. If a tiling in one of these families is parametrized with
an integer scale \(t\) and tile count \(t^2(3v^2-u^2)\) or
\(t^2(2v^2-u^2)\), then a squarefree tile count forces \(t=1\) and is
excluded by the corresponding obstruction. This is a consequence within
these two parametrized branches only; it does not classify all
squarefree tile counts or all composite tile counts in Erdős problem
634.

\clearpage
\begin{thebibliography}{9}

\bibitem{laczkovich}
M. Laczkovich, \emph{Tilings of triangles}, Discrete Mathematics
\textbf{140} (1995), 79--94.
\url{https://doi.org/10.1016/0012-365X(93)E0176-5}.

\bibitem{reptile}
S. L. Snover, C. Waiveris and J. K. Williams,
\emph{Rep-tiling for triangles}, Discrete Mathematics
\textbf{91} (1991), 193--200.
\url{https://doi.org/10.1016/0012-365X(91)90110-N}.

\bibitem{isosceles}
M. Beeson, \emph{Tilings of an Isosceles Triangle},
arXiv:1206.1974v7, 4 May 2026.
\url{https://arxiv.org/abs/1206.1974v7}.

\bibitem{group1}
M. Beeson, \emph{Triangle Tiling: The Case $3\alpha+2\beta=\pi$},
arXiv:1206.2229v4, 25 September 2026.
\url{https://arxiv.org/abs/1206.2229v4}.

\bibitem{equilateral}
M. Beeson, \emph{Tiling an Equilateral Triangle},
arXiv:1812.07014v3, 28 May 2024.
\url{https://arxiv.org/abs/1812.07014v3}.

\bibitem{rationality}
M. Beeson and Y. X. Zhang, \emph{Rationality of certain triangle
tilings}, arXiv:2604.01314v1, 1 April 2026.
\url{https://arxiv.org/abs/2604.01314v1}.

\bibitem{prime}
M. Beeson, \emph{Tiling a triangle into a prime number of congruent
triangles}, arXiv:2607.23453v1, 26 July 2026.
\url{https://arxiv.org/abs/2607.23453v1}.
Only the specified classification statements and individual branch
computations are used; the global prime-count claim is not an input.

\bibitem{bonfioli}
V. Bonfioli, \emph{A signed-direction invariant for triangle tilings},
repository snapshot, commit
\texttt{4bb61193bafa4471f7ba3a35324ae0f70bd2177c};
Zenodo version 3.2.
\url{https://github.com/ElVec1o/erdos_634_proof/tree/4bb61193bafa4471f7ba3a35324ae0f70bd2177c}.
\url{https://doi.org/10.5281/zenodo.22721069}.
The isosceles $120^\circ$ argument is reproduced here; broad prime
claims in the repository metadata are not used.

\end{thebibliography}

\end{document}
